\section{Degree-separated coding: a ten-integer equation}
\label{sec:wave9-signed}

There is a completely fixed integer polynomial
\[
 F_{10}(t,g_0,h_0,k_0,l,w,x,y,z,m)
\]
of exact total degree $3\,327\,556$, evaluated by a division-free circuit of
$519$ binary arithmetic operations, for which ordinary mathematics
formalizable in ZFC proves
\[
 \neg\exists(t,g_0,h_0,k_0,l,w,x,y,z,m)\in\mathbb Z^{10}\;(F_{10}=0)
 \quad\Longleftrightarrow\quad\operatorname{Con}(\mathrm{ZFC}).
\]
This combines the preceding degree-separated source architecture with the incoming congruence-norm finalizer. Neither ingredient alone is claimed as new. The contribution is their proved compatibility, including the constructive residue restrictions, the smaller source layout, and removal of an irrelevant low mask bit together with the leading high bit of the binomial argument.
The inherited compiler calibration and Sun arithmetic bridge remain stated
mathematical dependencies. No consistency decision, global minimality claim,
or new end-to-end proof-assistant certificate is asserted.

The structural change is to split the extraction convolution by homogeneous
degree. A linear label may now collide with a quadratic label: their kernels
are separate summands. This relaxes the finite isolation problem rather than
merely finding a better placement for the previous problem. Elementary row
operations, one shared positivity shift, borrow-tolerant target padding,
and exact mask-population calibration complete the construction.

\subsection{Fixed source and three homogeneous kernels}

Use exactly the unshifted, strictly positive, 31-coordinate parent in
\texttt{data/parent\_quartic31.json}, with its complete original four-port
binding $(\beta_{\rm native},K_0,M_C,M_F)$. The four included compressed
payloads have bit lengths $(1064921,1070004,1064921,1064922)$; the last is
\texttt{MF\_source}, already shifted, and is not shifted again.
\texttt{INSTANCE.json} is the active binding. No package-52 transport row or
coefficient is spliced into this tuple. The pinned parent theorem is
\[
 \exists p\in\mathbb Z_{>0}^{31}\;r_1(p)=\cdots=r_{22}(p)=0
 \quad\Longleftrightarrow\quad\neg\operatorname{Con}(\mathrm{ZFC}).
\]
Define $\rho_{14}=-(r_{14}+r_2)$, $\rho_{15}=-(r_{15}+r_3)$, and $\rho_j=-r_j$ otherwise. The inverse row
operations preserve the common zero set over every coefficient tuple.
There are now 89 records: seven constant, 39 linear, and 43 quadratic.

In the parent coordinate order
\[
\begin{split}
 &(J,F,\alpha,z,f,h,i,o,j,s,w,\tau,\eta,\zeta,y,\gamma,\\
 &\hspace{22mm}N,E,P,c,a,X,D,C_2,L,F_2,V,g,I,B_{\rm aux},\Delta),
\end{split}
\]
choose the positive, distinct weights
\[
\begin{split}
 (w_1,\ldots,w_{31})={}&(104,76,79,60,108,8,99,55,36,54,98,11,7,10,44,32,\\
 &106,22,17,1,3,4,37,89,97,107,39,40,86,45,92).
\end{split}
\]
The row order and targets are
\[
\begin{split}
 &(8,15,20,9,4,19,7,12,13,17,16,10,5,21,1,3,11,6,14,18,22,2),\\
 &(34,229,274,319,524,533,573,767,820,1007,1053,1064,1089,1291,\\
 &\hspace{16mm}1320,1366,1381,1410,1617,1707,1798,1939).
\end{split}
\]
Let $q_j$ denote the target assigned to row $j$, and write
$\rho_j=\sum_\alpha a_{j,\alpha}p^\alpha$. For $d=|\alpha|$ put
\[
 e(\alpha)=\sum_i\alpha_iw_i,\qquad
 c_{j,\alpha}=
 \begin{cases}2a_{j,\alpha},&d=0,1,\\
 a_{j,\alpha}\prod_i\alpha_i!,&d=2.
 \end{cases}
\]
Define the fixed integer polynomials
\[
 D_d(Z)=\sum_{j,\,|\alpha|=d}c_{j,\alpha}Z^{q_j-e(\alpha)},\qquad
 U_p(Z)=\sum_i p_iZ^{w_i},\qquad
 G_p(Z)=D_2(Z)U_p(Z)^2+D_1(Z)U_p(Z)+D_0(Z).
\]
All reversed exponents are nonnegative. The exact finite certificate checks,
for every used degree-$d$ monomial $\alpha$ and every degree-$d$ monomial
$\gamma$ in 31 variables,
\[
 q_j-e(\alpha)+e(\gamma)=q_k
 \quad\Longleftrightarrow\quad j=k\text{ and }\alpha=\gamma.
\]
There are respectively $1,31,496$ expansion monomials in degrees $0,1,2$,
and 22,544 source/expansion interactions. There is no requirement to
separate a degree-one label from a degree-two label. Multinomial expansion
and the displayed coefficient weights therefore give the universal identities
\begin{equation}\label{eq:split-extraction}
 [Z^{q_j}]G_p(Z)=2\rho_j(p)\quad(1\le j\le22).
\end{equation}
The independent checker compares the entire extracted polynomials in all
31 coordinates and four ports, not just evaluations at test points.

Write $M=108$, $Q=1939$. The exact kernel degrees and convolution height are
\[
 \deg D_0=1939,\qquad\deg D_1=1842,\qquad\deg D_2=1738,\qquad
 H=\max_d(\deg D_d+dM)=1954.
\]
The leading coefficients of $(D_0,D_1,D_2)$ are respectively $(2,2,\beta_{\rm native})$, all positive. Every kernel separately has distinct exponents. Replacing each
port monomial in a coefficient by its strict bit bound and taking absolute
integer coefficients gives a maximum majorant of bit length $2134927$.
Consequently, with
\[
 C_0=2^{2134927},\qquad\kappa=2134940,\qquad\beta_S=2^\kappa,
\]
every kernel coefficient has absolute value less than $C_0$, and
$\beta_S>4096C_0$.

\subsection{Source polynomials and pre-binary bounds}

In ten unrestricted integer unknowns $(t,g_0,h_0,k_0,l,w,x,y,z,m)$ define first
\[f=2t+1,\quad b=6t+4,\quad g=bg_0,\quad h=bh_0,\quad k=b(3k_0+2).\]
The circuit variable names are $(f,g,h,k,l,w,x,y,z,m)$; its first four names denote the free coordinates $(t,g_0,h_0,k_0)$, not these computed registers. Then define
\begin{gather*}
 b=1+3f,\quad B=\beta_S b^2,\quad\ell=\sum_i B^{w_i},\quad U_c=\ell+g,\\
 G_c=D_2(B)U_c^2+D_1(B)U_c+D_0(B),\quad A_T=\sum_j B^{q_j},\\
 N_0=B^{109},\quad N_1=B^{1955},\quad N=N_0N_1=B^{2064},\\
 M_0=N_0-1-(b-1)\ell,\quad J_c=2G_c+A_T,\\
 S_c=g+J_cN_0,\quad T_c=M_0+(B-2)A_TN_0-2,\\
 X=(S_c+T_c)N+N-1-S_c,\\
 Y=\beta_S^{15}\,2^{\kappa-23}B^{2179}b
   =2^{2195\kappa-23}b^{4359},\qquad
 \mathcal C=bB^{108}=2^{108\kappa}b^{217}.
\end{gather*}
All powers have fixed nonnegative integer exponents. Every displayed
register is an abbreviation, never an additional quantified unknown.

Assume $f\ne0$, $b$ is a square, and $0\le g<2\mathcal C$.
Then $b\equiv1\pmod3$ implies $b\ge4$, and $B>4096C_0b^2$.
Distinct positive weights give $\ell<2B^{108}$, hence
\[
 B^{108}\le U_c<3bB^{108}.
\]
Every lower coefficient has absolute value below $C_0$. For a kernel of degree $d$ with positive integer leading coefficient $c$, the sum of absolute values of its lower evaluated terms is less than
\[
 C_0\frac{B^d-1}{B-1}<\tfrac12B^d.
\]
Consequently all three $D_j(B)>0$, and $G_c>0$ for $U_c\ge0$. This proves the pre-binary positivity without inferring it merely from positive inputs.

The kernel degree bounds also give
\[
 |G_c|<2C_0(1+3b+9b^2)B^H\le26C_0b^2B^H,\qquad
 2G_c<N_1/64.
\]
Since $Q=1939<H$, both $A_T$ and $(B-2)A_T$ are less than $N_1/64$.
It follows that
\[
 0<J_c<N_1/32,\qquad J_c+(B-2)A_T+2<N_1.
\]
The digits of $M_0$ are $B-1$, except that the 31 designated cells
have digit $B-b$. Therefore $0\le M_0<N_0$, and $g<N_0$ as well.
The preceding strict high-block bound proves
\begin{equation}\label{eq:split-tight-range}
 0\le S_c,\quad0\le T_c,\qquad S_c+T_c<N.
\end{equation}
All source weights are positive, so $M_0\ge B-1$ and $M_0-2\ge B-3>0$. Moreover $S_c\ge N_0$, because $J_c$ is a positive integer and $g\ge0$. Hence
\[
X=S_c(N-1)+T_cN+N-1>N_0(N-1)>3b.
\]
The reduced $Y$ still exceeds $\max(b,256)$. Because $\kappa$ is
even, $B\equiv N\equiv1\pmod3$, $M_0\equiv0$, and
$T_c\equiv-24$; hence $X\equiv T_c\equiv0\pmod3$.
These are all established before binary-base recovery.

\subsection{Borrow-tolerant extraction and exact mask population}

Suppose now that $b$ is a power of two. So are $B,N_0,N$. The low
mask uses $M_0-2$ instead of $M_0$. This clears precisely bit 1, since the zero-th base-$B$ cell of $M_0$ is $B-1$. After recovery of a binary square base, $b\ge4$ and $g=bg_0$ is divisible by four, so its bit 1 is zero already. Thus $g\mathbin{\&}(M_0-2)=0$ is equivalent to $g\mathbin{\&}M_0=0$, which is exactly
\[
 g=\sum_i n_iB^{w_i},\quad0\le n_i<b.
\]
For these digits put $p_i=n_i+1$, so $U_c=\sum_i p_iB^{w_i}$ encodes
positive parent coordinates without expanding a successor substitution
in every residual. Write $h_s=[Z^s]G_p(Z)$ and let $\epsilon_s$ be the
indicator of the 22 targets. All coefficients, including unsupported
collisions, satisfy
\[
 |h_s|<C_0\{(\sum_i p_i)^2+\sum_i p_i+1\}
 \le993C_0b^2<B/4.
\]
The signed expansion of $J_c$ is $\sum_{s=0}^{H}(2h_s+\epsilon_s)B^s$.
Euclidean digit normalization has incoming carries only in $\{-1,0\}$:
\[
 u_0=0,\quad u_{s+1}=\left\lfloor\frac{2h_s+\epsilon_s+u_s}{B}\right\rfloor,
 \quad d_s=2h_s+\epsilon_s+u_s-Bu_{s+1}.
\]
At a marked cell the mask $B-2$ admits precisely digits $0,1$.
If $h_s=0$, the unnormalized digit $1+u_s$ passes. If $h_s\ge1$,
it lies in $[2,B)$ and fails. If $h_s\le-1$, it lies strictly between
$-B+1$ and zero, and its normalized digit is at least two, so it fails.
Thus incoming borrows do not affect the zero test. Positivity of $J_c$
was proved independently, ensuring these are its ordinary digits.
Combining the disjoint low and high blocks with
\eqref{eq:split-extraction} proves
\[
 S_c\mathbin{\&}T_c=0
 \quad\Longleftrightarrow\quad
 g=\sum_i n_iB^{w_i},\ 0\le n_i<b,\quad r_1(n+1)=\cdots=r_{22}(n+1)=0.
\]

For $N=2^u$ and $S,T\ge0$ with $S+T<N$, let
$X_*=(S+T)N+N-1-S$. Its disjoint binary blocks give
\[
 s_2(X_*)=u+s_2(T)-\{s_2(S)+s_2(T)-s_2(S+T)\}.
\]
The braces count carries with multiplicity, and vanish precisely when
$S\mathbin{\&}T=0$. Legendre's formula therefore implies
\[
 2^{u+s_2(T)}\mid\binom{2X_*}{X_*}
 \quad\Longleftrightarrow\quad S\mathbin{\&}T=0.
\]
The stronger range \eqref{eq:split-tight-range} is essential.
Writing $b=2^k$, $B=2^v$, the actual mask population is
\[
 s_2(T_c)=(109+22)v-31k-23,\qquad u=2064v.
\]
Hence $u+s_2(T_c)=2195v-31k-23=\log_2Y$, exactly explaining the
fixed polynomial threshold used above.

Conversely, a positive parent zero is not the all-one vector: original
row five equals one there. For arbitrarily large $b=4^s$ exceeding all
its coordinates, encode $n_i=p_i-1$. Then $g\ge B>b$ and $g<\mathcal C$;
the latter follows from $(b-1)\ell<bB^{108}$. All marked tests pass and
$Y\mid\binom{2X}{X}$. This proves both directions of the entire source
contract needed by the inherited finalizer, including its lower bound on
forward witnesses and its pre-binary reverse bounds.

\subsection{The Sun bridge registers}\label{the-sun-bridge-registers}

Set \[
 L=lY,\quad U=2LX,\quad V=4gwY,\quad W_*=bw,
\] \[
 A=U(V+1),\quad B_*=2X+1,\quad C=B_*+(A-2)h,
\] \[
 P=U^2V,\quad K=X+1+k(P-2),\quad
 S=2A-5,\quad T=3W_*C-2(W_*^2-1),
\] \[
 Q=\bigl((P/2)^2-1\bigr)K^2+1,
\] \[
 \Delta=8\mathcal C^3gK^2-
 g^2\bigl(32(C-KL)^2\mathcal C^3+g^2K^2\bigr).
\] The notation \(P/2\) is only an abbreviation for the division-free
expression \((LX)UV\), as \(U=2LX\); \texttt{Phalf} is exactly that
expression in the circuit. Thus \(4Q=(P^2-4)K^2+4\), the old auxiliary
square radicand.

Put \(\xi=zx\), and form Sun's modified Lucas certificate \[
 D=(A^2-4)C^2+4,\qquad E=C^2D\xi,\qquad
 F=4(A^2-4)E^2+1,
\] \[
 G_*=2(A+2)(A-2)^2E^2-1,\qquad \Omega=16E,
\] \[
 H_*=C+(B_*-C)F+\Omega Fy,\qquad I_*=(G_*^2-1)H_*^2+1.
\] Finally set \[
 \rho=T+Sy,\quad \Sigma=\xi S^5,\quad
 \Theta=(\rho^3-\rho S^2-S^3)(\rho^2+23S^2).
\]

The external bridge input is precisely Sun (2021), Theorems 4.12--4.13,
specialized to \(p=P_{\rm theorem}=2,Q_{\rm theorem}=1\), with the
modified Lucas certificate and residue interpolation of Sun (2026),
Theorems 3.1 and 4.1 and Propositions 6.1--6.2. These give the following
interface.

\begin{itemize}
\item
  From a source witness with \(b\in2\uparrow\), \(0<g<\mathcal C\), and
  the central-binomial divisibility, choose positive \(l,w\) and
  integers \(h,k\) so that \(Q\) is a square, \(S\mid T\), \(\Delta>0\),
  \(C=u_{2X+1}(A)\), and \(K=u_{X+1}(P)\).
\item
  Conversely, if \(b\) is square, \(f^2l^2\Delta>0\), \(\xi\ne0\),
  \(4Q\) and \(DFI_*\) are squares, and \(S\mid T\), the bridge recovers
  \(b\in2\uparrow\), \(0<g<2\mathcal C\), and the binomial divisibility.
\item
  On the forward branch, \(\xi\) can be made a multiple of any
  prescribed positive integer, and \(y\) can be assigned any residue
  modulo \(\xi\) while retaining the square certificate \(DFI_*\).
\end{itemize}

The exact forward values of \(C,K\) are not an extra assumption: the
proof of Sun's 2021 Theorem 4.12 defines \(C=u_B(A,1)\) immediately
before equation (4.16), which defines \(K=u_{QX+1}(U^PV,1)\). The
recurrence used here is
\(u_0(a)=0,u_1(a)=1,u_{n+1}(a)=au_n(a)-u_{n-1}(a)\).

\subsection{Congruences valid for every integer
assignment}\label{congruences-valid-for-every-integer-assignment}

These identities do not assume positivity, a source solution, a Pell
equation, or a zero of the final polynomial.

First, \(B,N_0,N\) are multiples of \(b^2\). Therefore \[
 X+1\equiv-g\equiv0\pmod b.
\] Also \(X\) is odd, because \(b,g,B,N_0,N\) are even. Since
\(b\equiv\beta_S\equiv1\pmod3\), all powers of \(B\) equal 1 modulo 3.
There are 22 targets. Modulo 3 one has \(M_0=0\), \(T_c=-24=0\), and
hence \[
 X\equiv(S_c+T_c)+1-1-S_c=T_c=0\pmod3.
\]

The displayed formula for \(Y\) has a factor \(b\); hence \(2b\mid U\)
and \(2b\mid A\). Consequently \(b\mid P\) and \(3\mid P\). The
normalization \(k=b(3k_0+2)\) gives \[
 K\equiv0\pmod b,\qquad K\equiv1+2(-2)=0\pmod3.
\] It follows that \[
 Q\equiv1\pmod{b^2},\qquad Q\equiv1\pmod3,
 \qquad \gcd(b,Q)=1.
\] In particular \(Q\ne0\), \(Q\) is odd, and for \(r=bQ\) we have
\(r\equiv1\pmod3\), \(|r|\ge2\).

Since \(h=bh_0\), \[
 C+1=2(X+1)+(A-2)bh_0\equiv0\pmod{2b}.
\] Consequently \(4b\mid D\): both terms of \(D=A^2C^2-4(C^2-1)\) are
divisible by \(4b\).

\subsection{The equation and unconditional two-root
separation}\label{the-equation-and-unconditional-two-root-separation}

Define \[
 r=bQ,\quad a=r+(3\Delta-1)(rz)^2,
 \quad d=DFI_*,\quad Z=\Sigma m-\Theta.
\] The single equation is \[
 \boxed{F_{10}=\bigl(Z^2-\Sigma^2(a+d)\bigr)^2
                  -4\Sigma^4ad=0.}
\]

The crucial separation is small and unconditional. Namely, \[
 \frac ab=Q+(3\Delta-1)bQ^2z^2\quad\hbox{is odd},\qquad
 \frac db\quad\hbox{is even}.
\] Both displayed quotients are integers; they are not operations in the
emitted polynomial. Since \(b\ne0\), this proves \(a\ne d\) on every
assignment. No magnitude comparison or large separating coefficient is
required.

The following elementary lemma explains why this is enough. If integers
\(a\ne d\) and a rational number \(q\) satisfy \[
 (q^2-a-d)^2=4ad,
\] then \(q\ne0\), because at zero the difference is \((a-d)^2\). Define
\(u=(q^2+a-d)/(2q)\) and \(v=q-u\). Expanding gives \(u^2=a,v^2=d\). A
rational square root of an integer is integral, so both \(a,d\) are
integer squares and \(q=u+v\in\mathbb Z\). This proof includes negative
and zero candidate radicands without making any prior choice of real or
complex square roots.

Now \(S=2A-5\) is odd and never zero. The polynomial
\(\pi(v)=(v^3-v-1)(v^2+23)\) has no rational root, so
\(\Theta=S^5\pi(\rho/S)\ne0\). If \(\Sigma=0\), the final polynomial is
\(\Theta^4\ne0\). At any zero we therefore have \(\Sigma\ne0\), and in
particular \(\xi\ne0\), \(z\ne0\). Apply the lemma with \(q=Z/\Sigma\).
It yields \[
 a,d\in\operatorname{Sq}(\mathbb Z),\qquad \Sigma\mid\Theta.
\] The latter follows from \(q=m-\Theta/\Sigma\in\mathbb Z\).

The standard quintic merger then gives \(S\mid T\). For clarity: if for
some prime \(p\) one had \(v_p(T)<v_p(S)\), then \(v_p(\rho)=v_p(T)\).
The two factors of \(\Theta\) have valuations \(3v_p(T)\) and
\(2v_p(T)\), respectively. Thus \(v_p(\Theta)=5v_p(T)<5v_p(S)\),
contradicting \(S^5\mid\Theta\).

\subsection{The combined square and positivity
guard}\label{the-combined-square-and-positivity-guard}

Factor \[
 a=r\{1+(3\Delta-1)rz^2\}.
\] The two factors are coprime. Neither is zero: \(|r|\ge2\), \(z\ne0\),
and the second factor is 1 modulo \(r\). Since their product is a
square, each is a signed square and both have the same sign. But
\(r\equiv1\pmod3\) excludes a negative square. Hence \(r\) and the
bracket are positive squares.

Similarly \(r=bQ\), \(\gcd(b,Q)=1\), and \(b\equiv Q\equiv1\pmod3\).
Each factor is a signed square, and the same modulus excludes negative
squares. Thus \(b,Q\) are positive squares. Since \(b\) is even and
nonzero, \(b\ge4\), and therefore \(r\ge4\).

If \(\Delta\le0\), then \[
 1+(3\Delta-1)rz^2\le1-r\le-3,
\] contradicting positivity of the bracket. Therefore \(\Delta>0\).
Since \(\mathcal C=bB^{108}>0\), the formula for \(\Delta\) gives
\(K\ne0\) and \(g>0\). If \(l=0\), then \(U=P=0\) and \(Q=1-K^2\le0\),
contradicting its positive-square value. Thus \(l\ne0\), and
\(f=2t+1\ne0\) was automatic. We have recovered \(f^2l^2\Delta>0\),
exactly the old bridge hypothesis, without placing the factor \(f^2l^2\)
in the new guard.

Before applying the bridge, rewrite
\[
\Delta=gK^2(8\mathcal C^3-g^3)-32g^2(C-KL)^2\mathcal C^3>0.
\]
Since $g>0$ and $K\ne0$, this gives $g^3<8\mathcal C^3$, hence $0<g<2\mathcal C$. The pre-binary source bounds already proved now yield $X>3b$ and $Y>\max(b,256)$ independently of any bridge conclusion. All hypotheses of the reverse bridge now hold. It yields the binomial
source contract, which decodes a positive zero of the pinned parent.
This proves soundness for all signed integer assignments, including
degenerate values of the individual bridge variables.

\subsection{Completeness of the new congruence
restrictions}\label{completeness-of-the-new-congruence-restrictions}

Suppose the parent has a positive zero. Choose a sufficiently large
\(b=4^s\), \(s\ge1\), in the cofinal source construction. Then
\(t=(b-4)/6\in\mathbb Z\), and the encoded \(g\) is a multiple of
\(b^2\). In fact \[
 X+1\equiv-g\equiv0\pmod{b^2}.
\] Apply the constructive bridge. It gives the exact values
\(C=u_{2X+1}(A)\), \(K=u_{X+1}(P)\) and positive \(l,w\), as stated
above. We show its \(h,k\) already have the imposed forms.

The index \(X+1\) is even. The recurrence shows that \(u_{2j}(P)\) is
divisible by \(P\). Since \(2b\mid U\), we have \(4b^2\mid P\), and thus
\(b^2\mid K\). Since \(P\equiv0\pmod4\), \(v_2(P-2)=1\). In \[
 (P-2)k=K-X-1,
\] the right side is divisible by \(b^2\). As \(b\) is a power of two,
\(b^2/2\mid k\), hence \(b\mid k\). Also \(P\equiv X\equiv0\pmod3\) and
\(K\equiv0\pmod3\), so \(k\equiv2\pmod3\). Because \(b\equiv1\pmod3\),
the integer \(k/b\) is 2 modulo 3. It is exactly of the form \(3k_0+2\).

For the other restriction, the recurrence gives
\(u_{2j+1}(A)\equiv(-1)^j\pmod{A^2}\). Since \(X\) is odd, \(C+1\) is
divisible by \(A^2\), hence by \(4b^2\). Also \(B_*+1=2(X+1)\) is
divisible by \(2b^2\). Thus \(2b^2\mid C-B_*\). The constructive \(A\)
is a multiple of four, so \(v_2(A-2)=1\). Therefore \[
 h=(C-B_*)/(A-2)
\] is divisible by \(b^2\), in particular by \(b\); write \(h=bh_0\).
All normalization restrictions thus preserve completeness, with no
change to the constructive \(l,w,C,K\).

These witnesses have \(\Delta>0\), and \(b,Q\) positive squares. Set
\(r=bQ\), also a positive square. The positive integer
\(\delta=3\Delta-1\equiv2\pmod3\) is nonsquare. Hence \(\delta r\) is a
positive nonsquare. Pell's theorem gives integers \(u,z>0\) with \[
 u^2-\delta r z^2=1.
\] It follows that \(a=ru^2\) is a square.

It remains to use the two free certificate variables without losing this
chosen \(z\). Since \(S\mid T\), write \(q=T/S\). First apply the
residue interpolation theorem with prescribed divisor \(z\) to obtain
its positive \(\xi\). Then choose a residue class of \(y\) for which
\(\pi(q+y)\equiv0\pmod\xi\). The polynomial \(\pi\) has a root modulo
every positive integer, so this residue exists. Interpolation supplies
the square \(d=DFI_*\), and \(x=\xi/z\) is integral. Then
\(\Theta=S^5\pi(q+y)\) is divisible by \(\Sigma=\xi S^5\). Choose \[
 m=\Theta/\Sigma+\sqrt a+\sqrt d\in\mathbb Z.
\] The polynomial identity for \(F_{10}\) now gives \(F_{10}=0\).

For the local-root fact used here, the cubic \(v^3-v-1\) has
discriminant \(-23\). At an odd prime other than 23, either it has a
simple linear factor, or its irreducibility makes the discriminant a
square modulo the prime, so \(v^2+23\) has a simple root. At 23 the
cubic has simple root 3. At 2, \(-23\equiv1\pmod8\) is a square in
\(\mathbb Z_2\). Hensel lifting and the Chinese remainder theorem prove
the assertion. The cubic has no rational root by the rational-root
theorem, and neither does \(v^2+23\). These are exactly the inherited
quintic merger facts.
\subsection{Circuit and exact degree certificate}
The artifact \texttt{generated/sun10.json} is a division-free circuit with ten unrestricted integer unknowns and 519 live binary arithmetic gates: 330 multiplications, 163 additions and 26 subtractions. Fixed coefficient construction is excluded from this arithmetic metric. The decisive degree ledger is
\[
\begin{array}{c|rrrrrr}
R&X&Y&A&C&K&Q\\\hline
\deg R&8254&4359&16976&16978&29592&118364\\
R&\Delta&a&d&\Sigma&Z&F_{10}\\\hline
\deg R&68559&305291&1494014&84882&84900&3327556
\end{array}
\]
The leading term is $\Sigma^4d^2$. Exact full-diagonal leading residues at degree $3327556$, modulo $101,1009,1000003$, are respectively $19,250,375721$. Any one nonzero residue proves that this upper degree is attained.

The supplied checks establish all 22,544 source-isolation interactions, 58 local circuit identities, the factored Lucas identity, the two-root norm, actual coefficient bindings, gate liveness, and the degree certificate. Additional finite borrow, population, source-size, congruence, and Lucas-recurrence tests supplement the universal arguments; they do not replace the imported arithmetic or compiler theorems. Both ordinary Python and optimized Python use explicit failures rather than disabled assertions.

Against our previous degree-separated ten-integer circuit, the degree falls from $8469528$ to $3327556$ and the operation count from $533$ to $519$. Against the incoming congruence-norm circuit, the degree falls from $9152284$ to $3327556$ and the count from $586$ to $519$. These comparisons do not claim global optimality.

\paragraph{Provenance and dependency boundary.}
The new tail is adapted from \path{con_zfc_congruence_norms_2026-10-23}, received in commit \texttt{7331974b696693a8eed2\allowbreak aa5f0f64f1c19d9569f0}; its complete proof is retained under \texttt{upstream/}. The source encoding is the newly certified split108 placement with a global sign reversal of the two-shear source. The prior split116 source proof is preserved for provenance only; the formulas above and active layout supersede its weights and kernels. The complete original four-port tuple is retained. Sun's 2021 Theorems 4.12--4.13 and 2026 Theorems 3.1, 4.1 and Propositions 6.1--6.2 remain explicit mathematical dependencies; the pinned source compiler theorem remains a dependency as well. This is not a new proof-assistant verification of the composite endpoint.
